% =======================================================
% ФУНДАМЕНТАЛЬНЫЕ ТЕМЫ
% =======================================================

\section{Встроенные типы данных, представление чисел в памяти, строки}

% -------------------------------------------------------
\subsection{Фундаментальные типы}

В C++ есть набор \textbf{встроенных (фундаментальных) типов}, которые понимает
сам компилятор без подключения заголовков:

\begin{center}
\begin{tabular}{lll}
\toprule
\textbf{Категория} & \textbf{Типы} & \textbf{Назначение} \\
\midrule
Целочисленные & \texttt{char}, \texttt{short}, \texttt{int}, \texttt{long}, \texttt{long long} & целые числа \\
\addlinespace
Беззнаковые & \texttt{unsigned char}, \dots, \texttt{unsigned long long} & неотрицательные целые \\
\addlinespace
С плавающей точкой & \texttt{float}, \texttt{double}, \texttt{long double} & вещественные числа \\
\addlinespace
Логический & \texttt{bool} & \texttt{true} / \texttt{false} \\
\addlinespace
Символьный & \texttt{char}, \texttt{wchar\_t}, \texttt{char8\_t}, \texttt{char16\_t}, \texttt{char32\_t} & символы \\
\addlinespace
Пустой & \texttt{void} & отсутствие значения \\
\bottomrule
\end{tabular}
\end{center}

\subsubsection{Размеры типов не фиксированы стандартом}

Стандарт C++ гарантирует только \textbf{минимальные} размеры и соотношения, а не
точные значения в байтах:

\[
\texttt{sizeof(char)} = 1 \le \texttt{sizeof(short)} \le \texttt{sizeof(int)}
\le \texttt{sizeof(long)} \le \texttt{sizeof(long long)}
\]

Гарантируется: \texttt{char} $\ge 8$ бит, \texttt{short} и \texttt{int} $\ge 16$ бит,
\texttt{long} $\ge 32$ бит, \texttt{long long} $\ge 64$ бит. По стандарту
\texttt{sizeof(char) == 1} \textit{всегда} (это единица измерения размеров).

\begin{lstlisting}
#include <cstdint>
// Типы с фиксированной шириной (если нужен точный размер):
int8_t   a;   // ровно 8 бит
int32_t  b;   // ровно 32 бита
uint64_t c;   // ровно 64 бита, беззнаковый
intptr_t d;   // целое, в которое влезает указатель
\end{lstlisting}

\subsection{Представление целых чисел в памяти}

\subsubsection{Беззнаковые числа}

Беззнаковое число хранится как обычное двоичное представление. Для $n$ бит
диапазон значений: $[0,\ 2^n - 1]$.

\begin{lstlisting}[numbers=none]
uint8_t x = 200;   //  1100 1000  = 128 + 64 + 8 = 200
\end{lstlisting}

При переполнении беззнаковые числа работают \textbf{по модулю} $2^n$ — это
определённое поведение (wrap-around):

\begin{lstlisting}
uint8_t x = 255;
x = x + 1;   // 0 (256 mod 256), НЕ undefined behavior
\end{lstlisting}

\subsubsection{Знаковые числа: дополнительный код (two's complement)}

Отрицательные числа хранятся в \textbf{дополнительном коде}. Старший бит ---
знаковый. Чтобы получить представление $-x$: инвертируем все биты $x$ и
прибавляем 1.

\begin{lstlisting}[numbers=none]
int8_t  5  =  0000 0101
int8_t -5  =  1111 1011   (инверсия 5 -> 1111 1010, +1 -> 1111 1011)
\end{lstlisting}

Преимущество дополнительного кода: сложение знаковых и беззнаковых чисел
выполняется \textbf{одной и той же} аппаратной операцией, и существует ровно один
ноль. Диапазон для $n$ бит: $[-2^{n-1},\ 2^{n-1} - 1]$ (отрицательных значений на
одно больше).

\textbf{Важно:} переполнение \textit{знакового} целого --- это
undefined behavior (UB), в отличие от беззнакового.

\subsubsection{Порядок байтов (endianness)}

Многобайтовое число можно хранить от младшего байта к старшему (little-endian,
x86) или наоборот (big-endian, сеть):

\begin{lstlisting}[numbers=none]
uint32_t x = 0x12345678;
// little-endian в памяти:  78 56 34 12
// big-endian    в памяти:  12 34 56 78
\end{lstlisting}

\subsection{Числа с плавающей точкой (IEEE 754)}

\texttt{float} (32 бита) и \texttt{double} (64 бита) хранятся по стандарту
IEEE~754: знак, экспонента (смещённый порядок) и мантисса.

\[
\text{значение} = (-1)^{\text{знак}} \times 1.\text{мантисса} \times 2^{\text{экспонента} - \text{смещение}}
\]

\begin{center}
\begin{tabular}{lccc}
\toprule
\textbf{Тип} & \textbf{Знак} & \textbf{Экспонента} & \textbf{Мантисса} \\
\midrule
\texttt{float}  & 1 бит & 8 бит  & 23 бита \\
\texttt{double} & 1 бит & 11 бит & 52 бита \\
\bottomrule
\end{tabular}
\end{center}

\textbf{Следствия:}
\begin{itemize}[leftmargin=2em]
    \item Не все числа представимы точно: \texttt{0.1 + 0.2 != 0.3}.
    \item Сравнивать вещественные через \texttt{==} опасно --- нужно сравнивать
          с допуском (epsilon).
    \item Есть специальные значения: \texttt{+inf}, \texttt{-inf}, \texttt{NaN}
          (причём \texttt{NaN != NaN}).
\end{itemize}

\begin{lstlisting}
#include <cmath>
bool almost_equal(double a, double b, double eps = 1e-9) {
    return std::fabs(a - b) < eps;
}
\end{lstlisting}

\subsection{Строки}

В C++ есть две принципиально разные «строки».

\subsubsection{C-строки (\texttt{const char*})}

Массив символов, завершающийся нулевым байтом \texttt{'\textbackslash 0'}.
Длина не хранится --- её вычисляют, идя до нуля (\texttt{strlen}, $O(n)$).

\begin{lstlisting}
const char* s = "hello";  // в памяти: h e l l o \0  (6 байт)
\end{lstlisting}

Минусы: ручное управление памятью, нет информации о длине, легко допустить
переполнение буфера.

\subsubsection{\texttt{std::string}}

Класс из \texttt{<string>}, управляющий динамическим буфером. Хранит указатель на
данные, длину и вместимость (capacity), сам выделяет и освобождает память.

\begin{lstlisting}
#include <string>
std::string s = "hello";
s += " world";          // динамически растёт
std::cout << s.size();  // длина за O(1)
\end{lstlisting}

\textbf{Small String Optimization (SSO).} Короткие строки (обычно до 15--22
символов) хранятся \textbf{прямо внутри объекта} \texttt{std::string}, без
выделения кучи. Это резко ускоряет работу с короткими строками.

\subsubsection{\texttt{std::string\_view} (C++17)}

Невладеющая «ссылка» на последовательность символов: указатель + длина. Не
копирует данные, не выделяет память. Идеален для передачи строк в функции
только для чтения.

\begin{lstlisting}
#include <string_view>
size_t count_spaces(std::string_view sv) {  // без копирования
    return std::count(sv.begin(), sv.end(), ' ');
}
count_spaces("hello world");        // из литерала
std::string s = "...";
count_spaces(s);                    // из std::string
\end{lstlisting}

\textbf{Опасность:} \texttt{string\_view} не продлевает время жизни данных. Если
исходная строка уничтожена, \texttt{string\_view} становится висячим.


% =======================================================
\section{Преобразования типов: явные и неявные}

% -------------------------------------------------------
\subsection{Неявные преобразования (implicit conversions)}

Компилятор сам вставляет преобразование, когда тип выражения не совпадает с
ожидаемым. Это происходит при присваивании, передаче аргументов, в
арифметике и т.\,д.

\subsubsection{Целочисленные продвижения (integral promotion)}

Типы меньше \texttt{int} (\texttt{char}, \texttt{short}, \texttt{bool}) перед
арифметикой продвигаются до \texttt{int}:

\begin{lstlisting}
char a = 10, b = 20;
auto c = a + b;   // a и b -> int, тип c == int (а не char!)
\end{lstlisting}

\subsubsection{Стандартные арифметические преобразования}

Когда операнды разных типов, они приводятся к «более широкому» общему типу:

\begin{lstlisting}
int    i = 5;
double d = 2.0;
auto x = i + d;   // i -> double, тип результата double
\end{lstlisting}

\subsubsection{Опасность: знаковое и беззнаковое}

Если в выражении встречаются знаковый и беззнаковый типы одинаковой ширины,
знаковый приводится к беззнаковому --- источник коварных багов:

\begin{lstlisting}
unsigned int u = 1;
int i = -1;
if (i < u)  // i -> unsigned: -1 становится 4294967295
    std::cout << "i < u";
else
    std::cout << "i >= u";   // <-- выполнится ЭТА ветка!
\end{lstlisting}

\subsubsection{Сужающие преобразования (narrowing)}

Преобразование, при котором значение может не поместиться или потерять точность
(\texttt{double} $\to$ \texttt{int}, \texttt{int} $\to$ \texttt{char}). В обычном
коде допускается с потерей данных, но \textbf{запрещено} в списковой
инициализации \texttt{\{\}}:

\begin{lstlisting}
int x = 3.9;       // OK, x == 3 (дробная часть отброшена)
int y{3.9};        // ОШИБКА компиляции: narrowing запрещён в {}
\end{lstlisting}

\subsubsection{Пользовательские неявные преобразования}

Конструктор с одним аргументом и оператор приведения задают неявные
преобразования для пользовательских типов:

\begin{lstlisting}
struct Meters {
    double value;
    Meters(double v) : value(v) {}      // неявный конструктор
    operator double() const { return value; }  // неявное приведение
};
Meters m = 5.0;       // double -> Meters (через конструктор)
double d = m;         // Meters -> double (через operator double)
\end{lstlisting}

Чтобы запретить неявное преобразование через конструктор, используют
\texttt{explicit}:

\begin{lstlisting}
struct Meters {
    explicit Meters(double v);  // теперь Meters m = 5.0; -- ОШИБКА
};
Meters m(5.0);   // только явный вызов
\end{lstlisting}

\subsection{Явные преобразования (casts)}

C++ предоставляет четыре именованных оператора приведения вместо «всемогущего»
C-style каста \texttt{(T)x}. Они подробно разобраны в отдельном разделе
о кастах; кратко:

\begin{center}
\begin{tabular}{lp{9.5cm}}
\toprule
\textbf{Каст} & \textbf{Назначение} \\
\midrule
\texttt{static\_cast<T>}      & безопасные «логичные» преобразования (числа, иерархия вверх, \texttt{void*}) \\
\addlinespace
\texttt{const\_cast<T>}       & снятие/добавление \texttt{const}/\texttt{volatile} \\
\addlinespace
\texttt{reinterpret\_cast<T>} & побитовая переинтерпретация (указатели, целые) \\
\addlinespace
\texttt{dynamic\_cast<T>}     & безопасное приведение вниз по полиморфной иерархии (с проверкой в runtime) \\
\bottomrule
\end{tabular}
\end{center}

\begin{lstlisting}
double d = 3.7;
int i = static_cast<int>(d);   // явно: i == 3, намерение видно
\end{lstlisting}

Преимущества именованных кастов перед C-style: их легко искать в коде, они
выражают намерение и не делают «слишком много» случайно.


% =======================================================
\section{Структуры и объединения}

% -------------------------------------------------------
\subsection{Структуры (\texttt{struct})}

\texttt{struct} объединяет несколько полей в один тип. В C++ \texttt{struct} и
\texttt{class} --- почти одно и то же; разница только в \textbf{умолчательной
видимости}: у \texttt{struct} члены \texttt{public} по умолчанию, у
\texttt{class} --- \texttt{private}.

\begin{lstlisting}
struct Point {
    int x;
    int y;
};
Point p{3, 4};        // агрегатная инициализация
std::cout << p.x;     // доступ через точку
Point* pp = &p;
std::cout << pp->y;   // доступ через указатель
\end{lstlisting}

\subsubsection{Выравнивание и padding}

Поля структуры размещаются с учётом \textbf{выравнивания} (alignment): адрес поля
должен быть кратен его требованию выравнивания. Поэтому компилятор вставляет
«дыры» (padding):

\begin{lstlisting}
struct Bad {
    char  a;   // 1 байт
    // 3 байта padding
    int   b;   // 4 байта (адрес кратен 4)
    char  c;   // 1 байт
    // 3 байта padding (для кратности размера)
};  // sizeof(Bad) == 12, а не 6!

struct Good {
    int   b;   // 4
    char  a;   // 1
    char  c;   // 1
    // 2 байта padding
};  // sizeof(Good) == 8
\end{lstlisting}

\textbf{Правило:} располагайте поля от больших к меньшим, чтобы уменьшить padding.

\subsection{Объединения (\texttt{union})}

В \texttt{union} все поля \textbf{делят одну и ту же память}. Размер объединения
равен размеру наибольшего поля. В любой момент времени «активно» только одно
поле.

\begin{lstlisting}
union Value {
    int    i;
    float  f;
    char   bytes[4];
};
Value v;
v.i = 0x40490FDB;
std::cout << v.f;   // прочитали те же байты как float ~= 3.14159
\end{lstlisting}

\textbf{Опасность:} чтение неактивного поля --- технически UB (хотя на практике
часто используется для побитовой интерпретации; правильный способ ---
\texttt{std::bit\_cast} в C++20 или \texttt{memcpy}).

\subsubsection{Tagged union и \texttt{std::variant}}

Чтобы безопасно знать, какое поле активно, к \texttt{union} добавляют тег:

\begin{lstlisting}
struct Tagged {
    enum { INT, FLOAT } tag;
    union { int i; float f; };
};
\end{lstlisting}

В современном C++ вместо ручного tagged union используют \texttt{std::variant}
(C++17) --- типобезопасное объединение, которое само следит за активным типом.


% =======================================================
\section{Указатели}

% -------------------------------------------------------
\subsection{Что такое указатель}

\textbf{Указатель} --- это переменная, хранящая \textbf{адрес} другого объекта
в памяти. Тип указателя говорит, как интерпретировать байты по этому адресу.

\begin{lstlisting}
int x = 42;
int* p = &x;     // p хранит адрес x  (& -- взятие адреса)
std::cout << *p; // 42 -- разыменование (доступ к значению по адресу)
*p = 100;        // меняем x через указатель -> x == 100
\end{lstlisting}

Размер указателя одинаков для всех типов и равен ширине адреса платформы (8 байт
на x86-64).

\subsection{Нулевой указатель}

\texttt{nullptr} (C++11) --- специальное значение «никуда не указывает».
Разыменование \texttt{nullptr} --- UB (обычно падение).

\begin{lstlisting}
int* p = nullptr;
if (p) { /* не выполнится */ }   // указатель конвертируется в bool
\end{lstlisting}

Используйте \texttt{nullptr}, а не \texttt{NULL} или \texttt{0}: \texttt{nullptr}
имеет собственный тип \texttt{std::nullptr\_t} и корректно работает с перегрузкой
функций.

\subsection{Указатели и \texttt{const}}

Расположение \texttt{const} меняет смысл:

\begin{lstlisting}
const int* p;        // указатель на const int: нельзя менять *p, можно p
int* const p;        // const-указатель: нельзя менять p, можно *p
const int* const p;  // нельзя менять ни *p, ни p
\end{lstlisting}

Читать справа налево: «\texttt{p} is a const pointer to a const int».

\subsection{Арифметика указателей}

Прибавление к указателю $n$ сдвигает адрес на $n \times \texttt{sizeof(T)}$ байт,
а не на $n$ байт:

\begin{lstlisting}
int arr[5] = {10, 20, 30, 40, 50};
int* p = arr;
std::cout << *(p + 2);  // 30 (сдвиг на 2*sizeof(int) = 8 байт)
++p;                    // теперь p указывает на arr[1]
std::cout << p - arr;   // 1 (разность указателей в элементах)
\end{lstlisting}

\subsection{Указатель на void и на функцию}

\begin{lstlisting}
void* vp = &x;     // «бестиповый» указатель, нельзя разыменовать напрямую
int* ip = static_cast<int*>(vp);   // нужно вернуть тип

int (*fp)(int, int) = &add;        // указатель на функцию
int r = fp(2, 3);                  // вызов через указатель
\end{lstlisting}

\subsection{Висячие и невалидные указатели}

\begin{itemize}[leftmargin=2em]
    \item \textbf{Висячий (dangling)} --- указывает на уже освобождённую память
          (после \texttt{delete} или выхода объекта из области видимости).
    \item \textbf{Утечка (leak)} --- объект выделен \texttt{new}, но указатель
          потерян, и \texttt{delete} никогда не вызовется.
    \item \textbf{Двойное освобождение} --- \texttt{delete} дважды по одному
          адресу --- UB.
\end{itemize}

В современном C++ сырые \texttt{new}/\texttt{delete} заменяют \textbf{умными
указателями} (\texttt{std::unique\_ptr}, \texttt{std::shared\_ptr}), которые
освобождают память автоматически.


% =======================================================
\section{Массивы}

% -------------------------------------------------------
\subsection{Встроенные (C-style) массивы}

Непрерывный блок однотипных элементов фиксированного размера, известного на
этапе компиляции:

\begin{lstlisting}
int arr[5] = {1, 2, 3, 4, 5};
arr[0] = 10;                  // индексация с нуля
std::cout << sizeof(arr);     // 20 == 5 * sizeof(int)
std::cout << sizeof(arr)/sizeof(arr[0]);  // 5 -- число элементов
\end{lstlisting}

\textbf{Нет проверки границ:} \texttt{arr[10]} --- UB, компилятор не остановит.

\subsubsection{Распад массива в указатель (array decay)}

При передаче в функцию массив \textbf{распадается} (decay) в указатель на первый
элемент --- информация о размере теряется:

\begin{lstlisting}
void f(int a[]) {       // на самом деле int* a
    std::cout << sizeof(a);  // размер УКАЗАТЕЛЯ (8), а не массива!
}
int arr[5];
f(arr);  // передаётся &arr[0], размер потерян
\end{lstlisting}

Поэтому в C-стиле размер передают отдельным параметром: \texttt{f(int* a, size\_t n)}.

\subsection{Многомерные массивы}

\begin{lstlisting}
int m[2][3] = {{1,2,3}, {4,5,6}};
// в памяти лежат подряд (row-major): 1 2 3 4 5 6
std::cout << m[1][2];   // 6
\end{lstlisting}

\subsection{\texttt{std::array} (C++11)}

Обёртка над C-массивом фиксированного размера, но с интерфейсом контейнера: знает
свой размер, не распадается в указатель, поддерживает итераторы.

\begin{lstlisting}
#include <array>
std::array<int, 5> a = {1, 2, 3, 4, 5};
std::cout << a.size();        // 5 -- размер не теряется
for (int x : a) { /* ... */ } // работают range-based for и алгоритмы
a.at(10);                     // бросает std::out_of_range (с проверкой)
\end{lstlisting}

\subsection{\texttt{std::vector} --- динамический массив}

Когда размер неизвестен на этапе компиляции или должен меняться, используют
\texttt{std::vector} --- массив, выделяемый в куче и растущий по мере
необходимости (подробно --- в разделе о контейнерах).

\begin{lstlisting}
#include <vector>
std::vector<int> v = {1, 2, 3};
v.push_back(4);          // динамически растёт
std::cout << v.size();   // 4
\end{lstlisting}


% =======================================================
\section{Устройство памяти, сегменты}

% -------------------------------------------------------
\subsection{Сегменты памяти процесса}

Адресное пространство процесса делится на области (сегменты):

\begin{center}
\begin{tabular}{lp{9cm}}
\toprule
\textbf{Сегмент} & \textbf{Что хранит} \\
\midrule
\texttt{.text} (code)   & машинный код программы, только для чтения \\
\addlinespace
\texttt{.rodata}        & константы и строковые литералы (только чтение) \\
\addlinespace
\texttt{.data}          & инициализированные глобальные/статические переменные \\
\addlinespace
\texttt{.bss}           & неинициализированные глобальные/статические (обнуляются) \\
\addlinespace
\textbf{Куча (heap)}    & динамическая память (\texttt{new}/\texttt{malloc}), растёт вверх \\
\addlinespace
\textbf{Стек (stack)}   & локальные переменные, кадры вызовов, растёт вниз \\
\bottomrule
\end{tabular}
\end{center}

\begin{lstlisting}[numbers=none]
  Высокие адреса
  +------------------+
  |      Stack       |  <- растёт вниз (локальные переменные)
  |        |         |
  |        v         |
  |                  |
  |        ^         |
  |        |         |
  |      Heap        |  <- растёт вверх (new/malloc)
  +------------------+
  |   .bss / .data   |  <- глобальные/статические
  +------------------+
  |  .rodata / .text |  <- код и константы
  +------------------+
  Низкие адреса
\end{lstlisting}

\subsection{Стек vs Куча}

\begin{center}
\begin{tabular}{lp{4.5cm}p{4.5cm}}
\toprule
& \textbf{Стек} & \textbf{Куча} \\
\midrule
Управление  & автоматическое (компилятор) & ручное (\texttt{new}/\texttt{delete}) \\
\addlinespace
Скорость    & очень быстро (сдвиг указателя) & медленнее (поиск блока) \\
\addlinespace
Размер      & ограничен (обычно 1--8 МБ) & ограничен ОЗУ \\
\addlinespace
Время жизни & до конца области видимости & пока не вызовут \texttt{delete} \\
\addlinespace
Фрагментация & нет & возможна \\
\bottomrule
\end{tabular}
\end{center}

\begin{lstlisting}
void f() {
    int x = 5;                 // на стеке, исчезнет при выходе из f
    int* p = new int(5);       // объект в куче, указатель p на стеке
    delete p;                  // память кучи нужно вернуть вручную
}                              // x уничтожается автоматически
\end{lstlisting}

\subsection{Время жизни и классы хранения}

\begin{itemize}[leftmargin=2em]
    \item \textbf{Автоматическое} (локальные переменные) --- на стеке, живёт до
          конца блока.
    \item \textbf{Статическое} (\texttt{static}, глобальные) --- в \texttt{.data}/\texttt{.bss},
          живёт всё время работы программы.
    \item \textbf{Динамическое} (\texttt{new}) --- в куче, живёт до \texttt{delete}.
    \item \textbf{Thread-local} (\texttt{thread\_local}) --- по копии на каждый поток.
\end{itemize}

\subsection{Переполнение стека (stack overflow)}

Стек ограничен. Слишком глубокая рекурсия или огромный локальный массив
исчерпывают его:

\begin{lstlisting}
void recurse() { recurse(); }   // бесконечная рекурсия -> stack overflow
void big()     { int huge[10'000'000]; }  // слишком большой кадр стека
\end{lstlisting}


% =======================================================
\section{Стек вызова функций}

% -------------------------------------------------------
\subsection{Кадр стека (stack frame)}

При каждом вызове функции на стеке создаётся \textbf{кадр (frame)}, содержащий:

\begin{itemize}[leftmargin=2em]
    \item аргументы функции (часть может передаваться в регистрах);
    \item адрес возврата (куда вернуться после функции);
    \item сохранённый указатель кадра вызывающего;
    \item локальные переменные.
\end{itemize}

\begin{lstlisting}
int square(int n) {       // кадр square: n, локальные
    int r = n * n;
    return r;
}
int main() {
    int x = square(5);    // создаётся кадр square поверх кадра main
}
\end{lstlisting}

\subsubsection{Как растёт стек}

При вызове кадр \textbf{кладётся} на вершину стека, при возврате ---
\textbf{снимается}. Стек растёт в сторону младших адресов:

\begin{lstlisting}[numbers=none]
main() вызывает square(5):

  | кадр main   |          | кадр main   |          | кадр main   |
  +-------------+   --->   +-------------+   --->   +-------------+
                           | кадр square |
                           +-------------+          (square вернулась,
   до вызова                во время вызова           кадр снят)
\end{lstlisting}

\subsection{Указатели стека}

\begin{itemize}[leftmargin=2em]
    \item \texttt{SP} (stack pointer) --- указывает на вершину стека.
    \item \texttt{BP/FP} (base/frame pointer) --- база текущего кадра, по нему
          адресуются локальные переменные и аргументы.
\end{itemize}

\subsection{Соглашения о вызове (calling convention)}

Соглашение определяет, как передаются аргументы (регистры/стек), кто чистит стек,
в каком регистре возвращается результат. Примеры: System~V~AMD64 (Linux),
Microsoft~x64 (Windows), \texttt{cdecl}, \texttt{stdcall}. На x86-64 первые
несколько целочисленных аргументов обычно идут в регистрах.

\subsection{Раскрутка стека (stack unwinding)}

При выбросе исключения стек \textbf{раскручивается}: кадры снимаются один за
другим, и для каждого локального объекта вызывается деструктор. Это основа RAII
и обработки ошибок (подробно --- в разделе об исключениях).

\subsection{Хвостовая рекурсия}

Если рекурсивный вызов --- последнее действие функции, компилятор может
\textbf{переиспользовать} текущий кадр (tail-call optimization), превращая
рекурсию в цикл и избегая роста стека:

\begin{lstlisting}
int fact(int n, int acc = 1) {
    if (n <= 1) return acc;
    return fact(n - 1, acc * n);   // хвостовой вызов
}
\end{lstlisting}


% =======================================================
\section{Ссылки}

% -------------------------------------------------------
\subsection{Что такое ссылка}

\textbf{Ссылка (lvalue reference)} --- это \textbf{псевдоним} (другое имя) для уже
существующего объекта. В отличие от указателя, ссылку нельзя «перенаправить» и
нельзя сделать «нулевой».

\begin{lstlisting}
int x = 10;
int& ref = x;     // ref -- другое имя для x
ref = 20;         // меняем x через ссылку -> x == 20
int* p = &ref;    // &ref -- это адрес x
\end{lstlisting}

\subsubsection{Ссылка vs указатель}

\begin{center}
\begin{tabular}{lp{4.5cm}p{4.5cm}}
\toprule
& \textbf{Ссылка} & \textbf{Указатель} \\
\midrule
Инициализация & обязательна при объявлении & необязательна \\
\addlinespace
Перенаправление & невозможно & возможно \\
\addlinespace
Может быть «нулём» & нет & да (\texttt{nullptr}) \\
\addlinespace
Синтаксис доступа & как к обычной переменной & нужно \texttt{*} / \texttt{->} \\
\addlinespace
Арифметика & нет & да \\
\bottomrule
\end{tabular}
\end{center}

\subsection{Ссылки в параметрах функций}

Передача по ссылке избегает копирования и позволяет менять аргумент:

\begin{lstlisting}
void increment(int& x) { ++x; }       // меняет оригинал
void print(const std::string& s);     // не копирует, но не меняет
\end{lstlisting}

\textbf{Правило:} большие объекты, которые не меняются, передавайте по
\texttt{const T\&} --- без копии и без риска модификации.

\subsection{Висячие ссылки}

Нельзя возвращать ссылку на локальную переменную --- она уничтожится:

\begin{lstlisting}
int& bad() {
    int local = 5;
    return local;   // ОШИБКА: ссылка на уничтоженную переменную (UB)
}
\end{lstlisting}

\subsection{Константные ссылки и продление жизни}

\texttt{const}-ссылка может привязаться к временному объекту и \textbf{продлевает
его время жизни} до конца жизни ссылки:

\begin{lstlisting}
const std::string& r = std::string("temp");  // временный живёт, пока жив r
\end{lstlisting}

\subsection{Rvalue-ссылки (C++11)}

\texttt{T\&\&} --- ссылка на временный объект (rvalue). Лежит в основе
move-семантики и perfect forwarding (подробно --- в отдельном разделе):

\begin{lstlisting}
std::string&& rr = std::string("temp");  // привязка к rvalue
void sink(std::string&& s);              // забирает ресурсы временного
\end{lstlisting}


% =======================================================
\section{Перегрузка функций}

% -------------------------------------------------------
\subsection{Что это}

\textbf{Перегрузка (overloading)} --- несколько функций с \textbf{одним именем},
но разными списками параметров. Компилятор выбирает нужную по аргументам вызова.

\begin{lstlisting}
int    add(int a, int b)       { return a + b; }
double add(double a, double b) { return a + b; }
std::string add(const std::string& a, const std::string& b) { return a + b; }

add(1, 2);          // int-версия
add(1.5, 2.5);      // double-версия
add("a"s, "b"s);    // string-версия
\end{lstlisting}

\subsection{Что участвует в выборе}

Перегрузки различаются по:
\begin{itemize}[leftmargin=2em]
    \item числу параметров;
    \item типам параметров;
    \item \texttt{const}/\texttt{volatile} и ссылочности (\texttt{\&}/\texttt{\&\&}).
\end{itemize}

\textbf{Не различаются} только по \textbf{типу возвращаемого значения} --- это
ошибка:

\begin{lstlisting}
int    f(int);
double f(int);   // ОШИБКА: отличается только возвращаемый тип
\end{lstlisting}

\subsection{Разрешение перегрузки (overload resolution)}

Компилятор делает это в три шага:
\begin{enumerate}[leftmargin=2em]
    \item \textbf{Набор кандидатов} --- все функции с таким именем, видимые в
          точке вызова.
    \item \textbf{Жизнеспособные (viable)} --- те, что подходят по числу
          аргументов и для которых существуют преобразования.
    \item \textbf{Лучшая} --- та, что требует «наименьших» преобразований.
          Точное совпадение лучше продвижения, продвижение лучше стандартного
          преобразования, стандартное лучше пользовательского.
\end{enumerate}

Если две функции одинаково хороши --- \textbf{неоднозначность (ambiguous)} ---
ошибка компиляции:

\begin{lstlisting}
void f(long);
void f(double);
f(5);    // ОШИБКА: int -> long и int -> double равноправны
\end{lstlisting}

\subsection{Перегрузка по \texttt{const} и ссылочности}

\begin{lstlisting}
struct S {
    void g() &;        // вызывается на lvalue
    void g() &&;       // вызывается на rvalue (временном)
    int  get() const;  // на const-объекте
    int& get();        // на неконстантном
};
\end{lstlisting}

\subsection{Перегрузка vs переопределение vs сокрытие}

\begin{itemize}[leftmargin=2em]
    \item \textbf{Overloading} --- одно имя, разные сигнатуры, \textit{в одной}
          области видимости.
    \item \textbf{Overriding} --- замена \texttt{virtual}-функции в наследнике с
          \textit{той же} сигнатурой.
    \item \textbf{Hiding} --- имя в наследнике скрывает \textit{все} перегрузки
          из базы (лечится \texttt{using Base::name;}).
\end{itemize}

\subsection{Значения по умолчанию}

Не путать с перегрузкой: один и тот же набор параметров, но часть имеет значение
по умолчанию:

\begin{lstlisting}
void log(const std::string& msg, int level = 0);
log("hello");      // level == 0
log("err", 3);     // level == 3
\end{lstlisting}


% =======================================================
\section{Пространства имён (namespace)}

% -------------------------------------------------------
\subsection{Зачем нужны namespace}

\textbf{namespace} группирует имена и предотвращает \textbf{конфликты имён} между
разными библиотеками. Вся стандартная библиотека лежит в \texttt{std}.

\begin{lstlisting}
namespace math {
    double pi() { return 3.14159; }
    struct Vector { double x, y; };
}
double p = math::pi();      // доступ через ::
math::Vector v;
\end{lstlisting}

\subsection{Использование имён из namespace}

\begin{lstlisting}
using namespace std;        // вносит ВСЕ имена std (не рекомендуется глобально)
using std::cout;            // вносит только cout (предпочтительно)
\end{lstlisting}

\textbf{Совет:} не пишите \texttt{using namespace std;} в заголовках и в глобальной
области --- это возвращает проблему конфликтов имён.

\subsection{Вложенные namespace}

\begin{lstlisting}
namespace company::project::utils {   // C++17, компактная запись
    void helper();
}
// эквивалентно вложенным namespace company { namespace project { ... } }
\end{lstlisting}

\subsection{Анонимный namespace}

Имена в безымянном namespace видны \textbf{только в этом файле} (аналог
\texttt{static} для внутренней линковки):

\begin{lstlisting}
namespace {
    int internal_counter = 0;   // не виден из других .cpp
}
\end{lstlisting}

\subsection{Псевдонимы namespace}

\begin{lstlisting}
namespace fs = std::filesystem;   // короткий псевдоним
fs::path p = "/tmp";
\end{lstlisting}

\subsection{Argument-Dependent Lookup (ADL)}

При вызове функции компилятор ищет её \textbf{также} в namespace, где определены
типы аргументов. Поэтому работает \texttt{std::cout << x} без явного \texttt{std::}
перед \texttt{operator<<}:

\begin{lstlisting}
namespace lib { struct S{}; void process(S); }
lib::S s;
process(s);   // находит lib::process через ADL, хотя нет lib::
\end{lstlisting}


% =======================================================
\section{Компиляция: этапы, ошибки}

% -------------------------------------------------------
\subsection{Этапы сборки программы на C++}

Превращение исходника в исполняемый файл проходит четыре стадии:

\begin{center}
\begin{tabular}{lp{9.5cm}}
\toprule
\textbf{Этап} & \textbf{Что делает} \\
\midrule
1. Препроцессинг & раскрывает \texttt{\#include}, \texttt{\#define}, \texttt{\#ifdef}; результат --- единый текст (translation unit) \\
\addlinespace
2. Компиляция & переводит C++ в ассемблер, затем в объектный код (\texttt{.o}); проверяет синтаксис и типы \\
\addlinespace
3. Ассемблирование & переводит ассемблер в машинный код (объектный файл) \\
\addlinespace
4. Линковка & связывает объектные файлы и библиотеки в исполняемый файл, разрешает символы \\
\bottomrule
\end{tabular}
\end{center}

\begin{lstlisting}[numbers=none]
main.cpp --[препроцессор]--> main.i --[компилятор]--> main.s
        --[ассемблер]--> main.o --[линковщик]--> a.out
\end{lstlisting}

\begin{lstlisting}
g++ -E main.cpp -o main.i   # только препроцессинг
g++ -S main.cpp -o main.s   # до ассемблера
g++ -c main.cpp -o main.o   # до объектного файла (без линковки)
g++ main.cpp -o app         # всё сразу
\end{lstlisting}

\subsection{Единица трансляции и модель компиляции}

Каждый \texttt{.cpp} компилируется \textbf{независимо} в свою единицу трансляции.
Заголовки (\texttt{.h}) физически \textbf{вставляются} препроцессором в каждый
\texttt{.cpp}, который их включает.

\subsubsection{Защита от повторного включения}

Чтобы заголовок не вставился дважды в одну единицу трансляции:

\begin{lstlisting}
#pragma once          // современный способ
// или классический include guard:
#ifndef MY_HEADER_H
#define MY_HEADER_H
// ... содержимое ...
#endif
\end{lstlisting}

\subsection{Объявление vs определение (ODR)}

\begin{itemize}[leftmargin=2em]
    \item \textbf{Объявление (declaration)} --- говорит «такой символ существует»
          (можно много раз).
    \item \textbf{Определение (definition)} --- предоставляет тело/память
          (\textbf{ровно один раз} на всю программу --- One Definition Rule).
\end{itemize}

\begin{lstlisting}
// header.h
int add(int, int);          // объявление (прототип)
extern int counter;         // объявление глобальной переменной

// source.cpp
int add(int a, int b) { return a+b; }  // определение
int counter = 0;                       // определение
\end{lstlisting}

\subsection{Виды ошибок}

\begin{center}
\begin{tabular}{lp{8.5cm}}
\toprule
\textbf{Вид} & \textbf{Когда и пример} \\
\midrule
Ошибка препроцессора & нет файла в \texttt{\#include}, незакрытый \texttt{\#ifdef} \\
\addlinespace
Ошибка компиляции & синтаксис, неизвестный тип, несоответствие типов \\
\addlinespace
Ошибка линковки & \texttt{undefined reference} (объявлено, но не определено), \texttt{multiple definition} (нарушен ODR) \\
\addlinespace
Ошибка времени выполнения & segfault, деление на ноль, исключение --- ловится только при запуске \\
\bottomrule
\end{tabular}
\end{center}

\subsubsection{Типичная ошибка линковки}

\begin{lstlisting}
// foo.h
void foo();           // объявили
// main.cpp
#include "foo.h"
int main() { foo(); } // используем, но определения foo() нет нигде
// -> undefined reference to `foo()'  (ошибка ЛИНКОВКИ, не компиляции)
\end{lstlisting}

\subsection{Полезные флаги компилятора}

\begin{lstlisting}
g++ -std=c++20 -Wall -Wextra -Wpedantic main.cpp   # стандарт + предупреждения
g++ -g main.cpp            # отладочная информация (для gdb)
g++ -O2 main.cpp           # оптимизация
g++ -fsanitize=address main.cpp   # детектор ошибок памяти (ASan)
\end{lstlisting}
